\chapter{John Willard Milnor}\label{ch1962johnwillardmilnor:john-willard-milnor-fields-medal-1962-a-comprehensive-report}

\begin{center}\rule{0.5\linewidth}{0.5pt}\end{center}

\textbf{Prepared:} July 2026
\textbf{Subject:} John Willard Milnor (b. February 20, 1931, Orange, New Jersey, USA)
\textbf{Position:} Professor Emeritus, SUNY Stony Brook
\textbf{Doctoral Advisor:} Ralph H. Fox (Princeton University)
\textbf{Fields Medal Citation:} \emph{``Proved that a 7-dimensional sphere can
have several differential structures; this led to the creation of the field of
differential topology.''}

\begin{center}\rule{0.5\linewidth}{0.5pt}\end{center}

\begin{prereqs} Differential topology (smooth manifolds, tangent bundles),
homotopy theory (fundamental group, covering spaces), basic algebraic topology.
\end{prereqs}

\begin{objectives} Explain the discovery of exotic spheres; describe the
h-cobordism theorem and its significance; understand the Kervaire--Milnor
classification of homotopy spheres. \end{objectives}

\section{Problem History}

Before Milnor's work \cite{milnor1956}, it was widely believed that a given topological manifold
possessed at most one differentiable structure. The classical examples---Euclidean
spaces, spheres, projective spaces, Lie groups---all came with a ``natural''
smooth structure.

The question of whether smooth structures are unique was part of a broader puzzle
in the classification of manifolds. Whitney showed in 1944 that every topological
manifold of dimension $\leq 3$ admits a unique smooth structure.

The cobordism theory of Thom (Fields Medal 1958) provided an algebraic framework
for classifying manifolds up to cobordism, but the finer question of distinguishing
homeomorphic but non-diffeomorphic manifolds required new invariants.

A related problem was the Generalized Poincaré Conjecture: is every homotopy
$n$-sphere homeomorphic to $S^n$? For $n \geq 5$, Smale had recently proved this
using the h-cobordism theorem, but the \emph{smooth} Poincaré conjecture remained
open.

The h-cobordism theorem (Smale, 1961) states that a simply connected h-cobordism
between manifolds of dimension $\geq 5$ is trivial. This reduced classification to
understanding the group of homotopy spheres $\Theta_n$ under connected sum.

Milnor's doctoral thesis on link isotopy showed his gift for new geometric
constructions. His work on characteristic classes gave him the technical arsenal
to attack the exotic sphere problem.

Kervaire was independently studying smooth structures, showing that some
topological manifolds admit no smooth structure at all. This raised the question:
can the same topological manifold carry multiple smooth structures?

\section{How Was It Solved}

Milnor's construction begins with \textbf{3-sphere bundles over $S^4$}. Identifying
$\mathbb{R}^4$ with the quaternions $\mathbb{H}$, a map $f: S^3 \to SO(4)$ is
given by $f(u)v = u^h v u^j$ where $h + j = 1$. The total space $M_k$ is a
smooth 7-manifold homeomorphic to $S^7$ for all odd $k$.

To show $M_k$ is \textbf{not diffeomorphic} to $S^7$, Milnor defined an invariant
using Thom's cobordism theory. If $M^7 = \partial B^8$:

\[\lambda(M) = 2q(B) - \tau(B) \mod 7\]

where $q(B) = p_2[B]$ is a Pontryagin number and $\tau(B)$ is the signature.
Computing: $\lambda(M_k) \equiv k^2 - 1 \mod 7$. For $k \not\equiv \pm 1 \mod 7$,
$\lambda \neq 0$, proving $M_k$ is exotic.

With Kervaire, Milnor classified all homotopy spheres. The result:

\[\Theta_7 \cong \mathbb{Z}/28\]

so there are exactly 28 distinct smooth structures on $S^7$. The general formula:

\[|\Theta_n| = \frac{2^{2n-2}(2^{n-1}-1)B_n}{n} \cdot \alpha_n\]

involves Bernoulli numbers, revealing a deep connection between differential
topology and number theory.

The \textbf{h-cobordism theorem} states: if $W$ is a simply connected h-cobordism
between closed $n$-manifolds with $n \geq 5$, then $W \cong M_0 \times [0,1]$.
The proof uses Morse theory and the Whitney trick.

Milnor's book \emph{Morse Theory} (1963) provided the definitive treatment,
leading from Morse inequalities to Bott periodicity via Morse theory on loop
spaces.

\section{Impact of This Discovery}

The discovery of exotic spheres \emph{founded differential topology} as an
independent field.

The Kervaire--Milnor classification initiated surgery theory, developed by
Browder, Novikov, and Wall, becoming the primary tool for classifying
high-dimensional manifolds.

The h-cobordism theorem established the paradigm of using analytic techniques to
prove topological results.

Milnor's total curvature theorem proved $\int_K |\kappa|\, ds \geq 4\pi$ for any
smooth knot. The Fáry--Milnor theorem showed total curvature $< 4\pi$ implies
unknottedness.

The Milnor fibration gave the invariant $\mu(f) = \dim_\mathbb{C} \mathcal{O}_n /
J_f$, central to singularity theory and algebraic geometry.

Milnor K-theory, with symbols $\{a_1, \ldots, a_n\}$ modulo the Steinberg
relation, connected algebraic K-theory to Galois cohomology. The Milnor conjecture,
proved by Voevodsky, is a major achievement of 21st-century mathematics.

Exotic spheres connect to mathematical physics: the choice of smooth structure on
compactified dimensions affects string theory.

\section{Modern Views and Improvements}

The classification extends to spin, string, and $G_2$-structures, with
$\Theta_n^{\text{Spin}}$ playing roles in string compactifications.

Surgery theory has been refined into the surgery exact sequence. Its failure in
dimension 4 remains one of the deepest open problems.

The smooth Poincaré conjecture in dimension 4 is still open.

Hill--Hopkins--Ravenel (2016) proved the Kervaire invariant one problem in
dimensions $2^n - 2$ for $n \geq 8$.

Active areas include exotic smooth structures on 4-manifolds (uncountably many on
$\mathbb{R}^4$) and gauge-theoretic invariants of exotic spheres.

\begin{center}\rule{0.5\linewidth}{0.5pt}\end{center}

\section*{Summary}
\addcontentsline{toc}{section}{Summary}

Milnor discovered exotic smooth structures on $S^7$, proving that a topological
manifold can carry multiple inequivalent differentiable structures, and his work
on the h-cobordism theorem and the Kervaire--Milnor classification founded
differential topology as a major branch of modern mathematics.
